% Numerical regression fixtures: five ERA tables and selected source-language comparisons.
\begin{tabular}{rrrrrr}\toprule
$f$ & $b$ & $Q_{\rm ref}$ & $Q_-$ & $Q_+$ & $100W/Q_{\rm ref}$ (\%)\\\midrule
$-8$ & 2 & 2.000000 & 1.994738 & 2.005254 & 0.526 \\
$-10$ & 1.6 & 1.788854 & 1.784090 & 1.793607 & 0.532 \\
$-16$ & 1 & 1.414214 & 1.410311 & 1.418108 & 0.551 \\
\bottomrule\end{tabular}

\begin{tabular}{ccccccc}\toprule
$b$ & $k/b$ & $Q_{512}$ & $Q_-$ & $Q_+$ & $W_{\rm rel}(\%)$ & $\Delta(\%)$\\\midrule
% BEGIN GENERATED table_2d
0.05 & 0.10 & 0.221979 & 0.208143 & 0.235996 & 12.55 & 0.1586 \\
0.2 & 0.10 & 0.531549 & 0.504780 & 0.558604 & 10.13 & 0.0497 \\
0.5 & 0.10 & 0.895772 & 0.858355 & 0.933727 & 8.41 & 0.0034 \\
1.5 & 0.05 & 1.624760 & 1.576099 & 1.675522 & 6.12 & 0.0299 \\
3 & 0.05 & 2.368417 & 2.304805 & 2.435609 & 5.52 & 0.0378 \\
% END GENERATED table_2d
\bottomrule\end{tabular}

\begin{tabular}{ccccc}\toprule
Mesh & $g_b$ & Selected $k/b$ & $Q_h$ & $W_{\rm rel}(\%)$\\\midrule
$32\times16$ & 1.053e-02 & 0.20 & 0.896204 & 34.64 \\
$64\times32$ & 5.189e-03 & 0.20 & 0.895913 & 22.72 \\
$128\times64$ & 2.577e-03 & 0.20 & 0.895854 & 16.88 \\
$256\times128$ & 1.284e-03 & 0.10 & 0.895802 & 11.29 \\
$512\times256$ & 6.408e-04 & 0.10 & 0.895772 & 8.41 \\
$1024\times512$ & 3.201e-04 & 0.05 & 0.895757 & 5.63 \\
\bottomrule\end{tabular}

\begin{tabular}{lrrrrrr}\toprule
Bias $b$ & 0.005 & 0.37 & 1.99 & 2 & 2.01 & 3 \\\midrule
Exact $Q$ & 0.100000 & 0.860233 & 1.994994 & 2.000000 & 2.005000 & 2.500000 \\
$Q_-$ & 0.095245 & 0.838124 & 1.944158 & 1.949037 & 1.953904 & 2.424932 \\
$Q_+$ & 0.104959 & 0.881648 & 2.044852 & 2.050102 & 2.055351 & 2.575068 \\
$W$ & 0.009714 & 0.043524 & 0.100694 & 0.101064 & 0.101447 & 0.150136 \\
$|Q_{\rm mid}-Q|$ & 1.02e-04 & 3.46e-04 & 4.89e-04 & 4.30e-04 & 3.72e-04 & 1.27e-12 \\\midrule
$W_s$ & 5.003e-03 & 4.304e-02 & 1.005e-01 & 1.009e-01 & 1.012e-01 & 1.500e-01 \\
$W_p$ & 4.104e-03 & 4.725e-04 & 1.989e-04 & 1.981e-04 & 1.976e-04 & 1.356e-04 \\
$W_d$ & 6.074e-04 & 1.320e-05 & 2.098e-07 & $\approx0$ & $\approx0$ & $\approx0$ \\
$100W_s/W$ & 51.50 & 98.88 & 99.80 & 99.80 & 99.81 & 99.91 \\
$100\eta$ & 2.11e+00 & 1.59e+00 & 9.70e-01 & 8.52e-01 & 7.34e-01 & 1.69e-09 \\
\bottomrule\end{tabular}

\begin{tabular}{ccccc}\toprule
Mesh size $h$ & $R_h$ & $|R_h-R|$ & $J_h$ & $J_h-J_{\rm ref}$\\\midrule
$1/16$  & 0.428979 & 0.073495 & $-12.553002$ & 2.384679\\
$1/32$  & 0.466470 & 0.036004 & $-13.763608$ & 1.174072\\
$1/64$  & 0.484770 & 0.017704 & $-14.349226$ & 0.588454\\
$1/128$ & 0.493430 & 0.009045 & $-14.638121$ & 0.299560\\
\bottomrule\end{tabular}

三个载荷下共计算 36 个区间，逐一与文献参考值比较，均有 $Q_-<Q_{\rm ref}<Q_+$。在 $f=-8$ 时，中心偏压恰为转折值 $b=2$，左侧采样点处于部分耗尽分支，右侧采样点处于完全耗尽分支。尽管两侧电势遵循不同的解析表达式，由三点能量构造的区间仍包含转折点电荷 $Q_{\rm ref}=2$。在 $f=-10$ 时，换算后的自由边界 $X=4/\sqrt{5}$ 落在单元内部。固定 $k/b=0.01$，将网格从 32 个单元加密至 256 个单元，区间相对全宽由 $2.593\%$ 降至 $0.532\%$，四层网格的区间均包含文献参考电荷 $4/\sqrt{5}$。表\ref{tab:branches}给出细网格结果。这组文献基准从边界导数给出电荷，因而可以直接检验由体内能量推得的上下界。当偏压采样跨越分支转折，或自由边界位于单元内部时，计算所得的电荷下界和上界仍包含文献参考值。

另取 $b=0.005$ 和 $b=3$，分别考察低偏压下的部分耗尽情形和完全耗尽情形。由两种情形的一维解析解得到的接触电荷分别为 $0.1$ 和 $2.5$。在 $n=64$、$k/b=0.1$ 时，相应的计算区间为 $[0.086036,\,0.114380]$ 和 $[2.424729,\,2.575271]$，均包含各自的解析电荷。

$n=256$ 时，在所取的 11 个步长中，$k=0.0025$ 给出最窄区间 $[0.997721,1.002279]$。继续减小步长，区间反而变宽，与推论\ref{cor:width}中的 $1/k$ 放大项一致。在该步长下，电荷区间的全宽约为 $0.004558$，其中精确能量的偏压差分贡献约 $0.002500$，能量包络误差贡献约 $0.002058$。两项量级接近，说明此时偏压差分与能量包络误差均显著影响区间宽度。二维整边接触算例也可由精确能量计算这两类贡献。

在 $b=0.5$ 时，将网格加密至 $1024\times512$（表\ref{tab:mesh}）。六层网格的恢复电荷均约为 $0.896$，最后一次加密时，恢复值相对前一层仅变化 $0.001738\%$。此时 $k/b=0.05$ 给出最窄区间 $[0.870759,0.921227]$。以同一 $|Q_{512}|$ 归一化，区间全宽由 $8.41\%$ 降至 $5.63\%$，说明恢复值已趋于稳定时，加密网格仍能进一步收紧电荷的上下界。

固定 $b=0.5$ 和 $128\times64$ 网格，三个候选步长中，$k=0.1$ 给出最小全宽 $0.151220$。沿用原通量恢复方法，逐次搜索选取 $k/b\approx0.1414$，全宽降至 $0.143281$，缩小 $5.25\%$。若保持原步长及两侧原始能量，仅在同一 RT0 空间内优化中心通量，全宽降至 $0.151049$，缩小 $0.11\%$。在这一偏压和网格下，调整步长对收紧上下界更有效。

表\ref{tab:fullcontact}显示，在低偏压 $b=0.005$ 时，能量包络误差占电荷区间全宽的 $48.50\%$。当 $b=0.37$ 时，即使用精确能量计算，两侧偏压与中心偏压之间的差分仍会产生区间宽度，这一部分占全宽的 $98.88\%$。图\ref{fig:fullmesh}给出了固定步长 $k=0.037$ 时的网格加密结果：区间全宽由 $0.051146$ 降至 $0.043068$，已接近用精确能量计算得到的宽度 $0.043039$。因此，在这一步长下，单靠加密网格能够进一步缩小的宽度已很有限。
